% Shared print layout. The web's provisional design and MathJax are independent.
\usepackage[utf8]{inputenc}
\usepackage[T1]{fontenc}
\usepackage[spanish,es-nodecimaldot]{babel}
\usepackage{lmodern}
\usepackage{amsmath,amssymb}
\usepackage{graphicx}
\usepackage{geometry}
\geometry{a4paper,left=20mm,right=20mm,top=25mm,bottom=22mm,headheight=16pt,headsep=8mm}
\usepackage{fancyhdr,lastpage}
\usepackage{enumitem}
\usepackage{xcolor}
\usepackage{pgfplots}
\pgfplotsset{compat=1.18}
\usepgfplotslibrary{fillbetween}
\usepackage[hidelinks,pdfencoding=auto]{hyperref}
\hypersetup{pdflang=es-ES,pdfauthor={},pdfsubject={Matemáticas Aplicadas a las Ciencias Sociales II}}
\pagestyle{fancy}
\fancyhf{}
\fancyhead[L]{\small Matemáticas Aplicadas a las Ciencias Sociales II}
\fancyhead[R]{\small 2.\textsuperscript{o} Bachillerato}
\fancyfoot[L]{\footnotesize Tema 3 --- Integrales | Relación guiada --- Edición 1.0}
\fancyfoot[R]{\footnotesize Página \thepage\ de \pageref{LastPage}}
\renewcommand{\headrulewidth}{0.4pt}
\setlength{\parindent}{0pt}
\setlength{\parskip}{0.5em}
\setlist{itemsep=0.3em,topsep=0.3em}
\newcommand{\apartado}[1]{\subsection*{#1}}
\pgfplotsset{
  revision/.style={
    width=0.94\linewidth,height=7.5cm,
    axis lines=middle,grid=major,grid style={gray!18},
    tick label style={font=\small},label style={font=\small},
    xlabel={$x$},ylabel={$y$},samples=160,
    legend style={at={(0.5,-0.16)},anchor=north,draw=none,font=\small},
    clip=true
  }
}
% Relación guiada: huecos para completar, cajas de comprobación y tablas.
\usepackage{array}
\renewcommand{\arraystretch}{1.55}
\newcommand{\hueco}[1][3cm]{\underline{\makebox[#1]{}}}
\newcommand{\dd}{\,dx}
\newcommand{\comprueba}{\par\smallskip\noindent\textit{Comprueba derivando:} \hrulefill\par}
\newcommand{\resultado}{\par\smallskip\noindent\textbf{Resultado:} \hrulefill\par}
\newcommand{\lineas}[1]{\par\vspace{2pt}\foreach \i in {1,...,#1}{\noindent\hrulefill\par\vspace{6pt}}}
\newcounter{ejercicio}
\newcommand{\ejercicio}[2]{\stepcounter{ejercicio}\par\medskip\noindent\textbf{#1\theejercicio.} #2\par}
\tolerance=2000\emergencystretch=4em
